%! AP Statistics — Unit 2, Lesson 2.9 — Group Activity — Answer Key
\documentclass[10pt]{article}
\usepackage{apstats-article}
\usepackage{apstats-key}

\begin{document}

\pageheader{Unit 2, Lesson 2.9}{Group Activity --- Answer Key}
\namepartnerperiod

\vspace{0.18in}

% ── SCENARIO ─────────────────────────────────────────────────────────────────
\begin{scenariobox}[Context]{sky}
\small
Each tier works with a real-world dataset that has a nonlinear relationship. Your group will
diagnose the nonlinearity, apply a log transformation, fit a regression to the transformed
data, and back-transform to make a prediction. Show all work and write answers in context.
\end{scenariobox}

\vspace{0.12in}

% ── TIER R ────────────────────────────────────────────────────────────────────
\begin{tcolorbox}[colback=white, colframe=black!40, boxrule=0.8pt, arc=2mm,
    left=4mm, right=4mm, top=3mm, bottom=3mm,
    title={\textbf{Tier R --- Remediate}}, fonttitle=\bfseries, halign title=center]
\small

\textbf{Context: Smartphone Resale Value}

A researcher recorded the age of used smartphones (years, $x$) and their resale value
(\$, $y$). The scatterplot curves downward (exponential decay). A $\log_{10}$ transformation
of $y$ produces a linear scatterplot. Technology output for the regression of $\log y$ on $x$:
\[
  \widehat{\log(\text{value})} = 2.80 - 0.18x \qquad r^2 = 0.94
\]

\vspace{8pt}

\begin{enumerate}[leftmargin=1.6em, itemsep=0.18in]

  \item Which variable was transformed? \ans{$y$ (resale value)} \quad
        What type of model is suggested? \ans{exponential}

  \item Write the back-transformed equation for predicting the original resale value.
        Show work.\\[8pt]
        \ansline{$\hat{y} = 10^{2.80 - 0.18x}$\enspace (or equivalently
        $\widehat{\text{value}} = 10^{2.80}\cdot(10^{-0.18})^x \approx 631 \cdot 0.661^x$)}

  \item Predict the resale value of a smartphone that is \textbf{3 years old}. Show all work.\\[8pt]
        \ansline{$\widehat{\log(\text{value})} = 2.80 - 0.18(3) = 2.26$;\enspace
        $\hat{y} = 10^{2.26} \approx \$182$}

  \item The residual plot for the \textit{original} linear fit showed a curved pattern.
        What does it mean that $r^2 = 0.94$ for the \textit{transformed} model?\\[8pt]
        \ansline{$r^2 = 0.94$ means the linear model for $\log y$ on $x$ explains 94\% of the
        variation in $\log(\text{value})$. Combined with a random residual plot, this indicates
        the exponential model fits the data well.}

\end{enumerate}
\end{tcolorbox}

\vspace{0.12in}

% ── TIER A ────────────────────────────────────────────────────────────────────
\begin{tcolorbox}[colback=white, colframe=black!40, boxrule=0.8pt, arc=2mm,
    left=4mm, right=4mm, top=3mm, bottom=3mm,
    title={\textbf{Tier A --- Approaching Proficiency}}, fonttitle=\bfseries, halign title=center]
\small

\textbf{Context: Kepler's Third Law (Planetary Orbits)}

An astronomer recorded orbital period ($y$, Earth years) and mean orbital radius ($x$,
AU) for eight planets. The scatterplot curves upward. After applying $\log_{10}$ to
\textbf{both} $x$ and $y$, the new scatterplot is nearly perfectly linear. Technology output:
\[
  \widehat{\log(\text{period})} = 0.001 + 1.499\,\log(\text{radius}) \qquad r^2 = 0.9999
\]

\vspace{8pt}

\begin{enumerate}[leftmargin=1.6em, itemsep=0.18in]

  \item Which variable(s) were transformed? \ans{both $x$ (radius) and $y$ (period)} \quad
        What type of model? \ans{power}

  \item Write the regression equation in transformed form (using variable names).\\[6pt]
        \ansline{$\widehat{\log(\text{period})} = 0.001 + 1.499\,\log(\text{radius})$}

  \item Write the back-transformed power-model equation. Show work.\\[8pt]
        \ansline{$\hat{y} = 10^{0.001 + 1.499\,\log x} = 10^{0.001}\cdot x^{1.499}
        \approx 1.002\,x^{1.499} \approx x^{1.499}$}

  \item Predict the orbital period of a planet with a mean radius of \textbf{5.2 AU}.
        Show all work and give the answer in Earth years.\\[8pt]
        \ansline{$\widehat{\log(\text{period})} = 0.001 + 1.499\,\log(5.2)
        = 0.001 + 1.499(0.716) = 1.074$;\enspace
        $\hat{y} = 10^{1.074} \approx 11.9$ years}

  \item The theoretical exponent in Kepler's Third Law is exactly 1.5. Does the data
        support this? Explain using the regression output.\\[8pt]
        \ansline{Yes. The estimated exponent is $b = 1.499$, which is extremely close to the
        theoretical value of 1.5. With $r^2 = 0.9999$, the power model fits the data almost
        perfectly, strongly supporting Kepler's Third Law.}

\end{enumerate}
\end{tcolorbox}

\vspace{0.12in}

% ── TIER E ────────────────────────────────────────────────────────────────────
\begin{tcolorbox}[colback=white, colframe=black!40, boxrule=0.8pt, arc=2mm,
    left=4mm, right=4mm, top=3mm, bottom=3mm,
    title={\textbf{Tier E --- Extension}}, fonttitle=\bfseries, halign title=center]
\small

\textbf{Context: Comparing Two Transformations}

A biologist studying fish length ($y$, cm) vs.\ age ($x$, years) has summary statistics for
two candidate models:

\vspace{6pt}
\renewcommand{\arraystretch}{1.4}
\begin{tabularx}{\linewidth}{@{} l X X @{}}
\rowcolor{sky} \textbf{Model} & \textbf{Regression equation} & \textbf{$r^2$} \\
Linear (no transformation) & $\hat{y} = 8.2 + 3.4x$ & 0.78 \\
Transformed ($\log y$ on $x$) & $\widehat{\log y} = 0.92 + 0.11x$ & 0.97 \\
\end{tabularx}

\vspace{10pt}

\begin{enumerate}[leftmargin=1.6em, itemsep=0.18in]

  \item Which model is more appropriate? State \textbf{two} pieces of evidence.\\[8pt]
        \ansline{The transformed (exponential) model is more appropriate. (1) Its $r^2 = 0.97$
        is substantially higher than the linear model's $r^2 = 0.78$, indicating a much better
        fit. (2) The transformed model should have a random residual plot, while the linear
        model's residual plot likely shows a curved pattern.}

  \item Write the back-transformed exponential equation for the better model. Show work.\\[8pt]
        \ansline{$\hat{y} = 10^{0.92 + 0.11x} = 10^{0.92}\cdot(10^{0.11})^x \approx 8.32\cdot 1.288^x$}

  \item Predict the length of a fish that is \textbf{7 years old} using the better model.
        Show all steps including the back-transformation.\\[8pt]
        \ansline{$\widehat{\log y} = 0.92 + 0.11(7) = 0.92 + 0.77 = 1.69$;\enspace
        $\hat{y} = 10^{1.69} \approx 49.0$ cm}

  \item The linear model predicts a length of $8.2 + 3.4(7) = 32.0$ cm.
        The transformed model predicts a different value. Explain why the two models
        give different predictions, and which you trust more.\\[8pt]
        \ansline{The linear model assumes length increases by the same amount each year (additive).
        The exponential model assumes proportional growth (multiplicative). Because $r^2$ is
        much higher and the residual plot of the transformed model should be random, the
        transformed model's prediction ($\approx 49$ cm) is more trustworthy.}

  \item A colleague proposes trying $\log y$ vs.\ $\log x$ instead (power model).
        Explain what additional check you would perform to decide between the exponential
        and power models, and what the residual plot would look like for the better choice.\\[8pt]
        \ansline{Fit both models and examine the residual plots of each. The better model will
        have a residual plot that shows random scatter with no systematic pattern. Also compare
        $r^2$: the model with the higher $r^2$ \textit{and} a random residual plot is preferred.}

\end{enumerate}
\end{tcolorbox}

\end{document}
